\chapter{ML-DSA-OSH Integration Guide}

\section*{Purpose}
This chapter defines the integration seam between the stable wrapper contract and the imported ML-DSA-OSH hardware implementation.

\section{Why the Engine Adapter Exists}
The AXI-Lite wrapper is the PS-visible contract boundary and should remain stable across hardware implementation changes. Directly coupling the wrapper to ML-DSA-OSH-specific stream ordering, key segmentation, or completion details would make software, documentation, and verification collateral vulnerable to avoidable churn. The engine adapter isolates that risk.

\section{Intended Boundary}
The wrapper owns register decode, digest staging, control semantics, status visibility, and signature readback windows. The adapter owns the downstream handshake for a signing operation. The project-owned ML-DSA-OSH shim then translates that stable adapter contract into the actual imported upstream sign-stream protocol.

\section{Current Real Attachment Path}
The current real-core path is:
\begin{center}
  `wrapper -> mldsa\_engine\_adapter -> mldsa\_osh\_shim -> combined\_top`
\end{center}

The chosen upstream integration point is `combined\_top` from the imported `ref\_combined/src/combined\_top.v` design, used in signing mode with the ML-DSA-87 security-level setting.

\section{Current Adapter Modes}
The adapter currently supports:
\begin{itemize}[leftmargin=*]
  \item `STUB`, which reproduces the existing fake result path used by software and contract tests
  \item `CORE\_PLACEHOLDER`, which keeps the adapter seam testable even without the imported core
  \item `MLDSA\_OSH`, which attaches the real imported sign path through project-owned shim logic
\end{itemize}

\section{Responsibilities of the Project-Owned Shim}
\begin{itemize}[leftmargin=*]
  \item convert the appliance digest input into the current PoC message representation expected by the upstream sign path
  \item partition the static PoC secret key according to the inspected upstream layout
  \item sequence the upstream streaming input transaction in the real inspected order
  \item capture the upstream output stream
  \item reorder upstream signature components into the wrapper-visible buffer layout
\end{itemize}

\section{Risks of Direct Coupling}
Direct coupling of the wrapper to ML-DSA-OSH internals would risk register-map churn, wrapper FSM churn, and avoidable software regressions when the core integration evolves. It would also blur ownership boundaries between vendored third-party source and project-owned appliance logic.

\section{Verification Strategy}
Verification should proceed in this order:
\begin{enumerate}[leftmargin=*]
  \item preserve wrapper-level testbench coverage for register behavior and control sequencing
  \item retain deterministic `STUB` mode as the local regression baseline
  \item validate the `MLDSA\_OSH` fallback error path in current Verilog-only local simulation
  \item add mixed-language simulation or synthesis-backed validation for the imported sign path
  \item compare signature length and output ordering against upstream expectations and known-good vectors
  \item keep software tests unchanged unless an explicit interface decision is approved and documented
\end{enumerate}

\section{Current Non-Goals}
This guide does not define production key management, final throughput tuning, or end-to-end board-level bring-up procedures. The current repository state integrates the real imported source behind the adapter boundary but does not yet claim full locally simulated cryptographic correctness.